\section{Introduction}

Reconstruction on a tree asks whether spins observed far from the root
retain information about its spin. For irregular trees, nonlinear
capacities can distinguish models with the same exponential growth rate.
Pemantle and Peres \cite{PP} developed capacity criteria and concave
recursions for critical Ising models on trees. The dependence of
reconstruction on the number of states is studied, in particular, by
Sly \cite{Sly} and Mossel, Sly, and Sohn \cite{MSS}.

For the symmetric four-state model, the fixed-channel capacity theorem
of \cite{Capacity} uses a moment-saving inequality for a preserved class
of posterior laws from \cite{Threshold}. We take this inequality and its
closure properties as an explicit assumption. Starting from that input,
we prove a capacity criterion for variable edge channels on every locally
finite tree. The main issue is uniformity as channel parameters approach
one: estimates obtained by bounding posterior denominators away from
zero no longer suffice. A bound using the posterior numerator instead
removes this restriction. A second estimate uses flow splitting to
control the accumulated information loss, including along long chains
with nearly perfect transmission.

\subsection{The model and capacity}
Let $T$ be a deterministic rooted locally finite tree with root $\rho$.
Every vertex has finitely many children, but their numbers need not have
a common bound. Write $L_n=\{v:|v|=n\}$ and $\ch(v)$ for the children
of $v$. The tree may be finite and may have dead ends. Give the root a
uniform spin in $[4]=\{1,2,3,4\}$. Each edge $e=(v,w)$ transmits its
parent spin through the channel
\begin{equation}\label{eq:channel}
 P_e(j\mid i)=\lambda_e\mathbf 1_{\{i=j\}}+\frac{1-\lambda_e}{4},
 \qquad 0<\lambda_e\leq1.
\end{equation}
The edge randomness is independent, and the tree and all parameters are
known to the observer. Set $\unif=(1/4,1/4,1/4,1/4)$ and
\begin{align}
 p^{(n)}&=\Prob(\sigma_\rho\in\cdot\mid\sigma_{L_n}),\notag\\
 a_n&=\E\TV(p^{(n)},\unif),\qquad
 m_n=4\E\lVert p^{(n)}-\unif\rVert_2^2.\label{eq:information}
\end{align}
An empty observation has posterior $\unif$. Both sequences decrease by
data processing, so their limits $a_\infty,m_\infty$ exist.
Reconstruction means $a_\infty>0$.

Put
\[
 r_e=\lambda_e^2,\qquad R_\rho=1,\qquad
 R_w=\prod_{e\in[\rho,w]}r_e.
\]
A flow $\theta$ is a nonnegative assignment to edges, conserved at each
nonroot vertex, with finite mass
$|\theta|=\sum_{w\in\ch(\rho)}\theta(\rho,w)$.
Conservation forces zero incoming flow at a leaf. Let $\partial T$ be
the infinite rays from $\rho$, and define
\begin{align}
 V(\theta)&=\sup_{\xi\in\partial T}
  \sum_{e=(v,w)\in\xi}(1-r_e^2)
             \left(\frac{\theta(e)}{R_w}\right)^2,\label{eq:potential}\\
 C(T)&=\sup\{|\theta|:V(\theta)\leq1\}.\label{eq:capacity}
\end{align}
The supremum over no rays in \eqref{eq:potential} is zero. The zero
flow is allowed, and $C(T)$ may be infinite. This is an $L^3$ nonlinear
capacity in the convention where the path potential has exponent $p-1$;
the corresponding edge resistances are
$\sqrt{1-r_e^2}/R_w$. When all channels have one parameter
$0<\lambda<1$, its positivity agrees with the capacity in
\cite{Capacity}.

\begin{theorem}\label{thm:main}
Assume the posterior-law inequality and closure properties in
Assumption~\ref{ass:input}. For the model \eqref{eq:channel} on every
deterministic rooted locally finite tree,
\[
 a_\infty>0\quad\Longleftrightarrow\quad C(T)>0.
\]
More precisely, put
\[
 K=\frac1{3000},\qquad \kappa=\frac{2K}{(1+9K)^3}.
\]
Then
\begin{equation}\label{eq:main-bounds}
 \frac1{20000}\frac{C(T)}{\sqrt{1+C(T)^2}}
 \leq m_\infty
 \leq\min\left\{3,\frac{C(T)}{\sqrt\kappa}\right\}.
\end{equation}
The fraction on the left is interpreted as one if $C(T)=\infty$.
All constants are independent of the degrees and channel parameters.
\end{theorem}

The equivalence follows from these bounds and
\begin{equation}\label{eq:tv}
 \frac{m_n}{8}\leq a_n\leq\frac{\sqrt{m_n}}2.
\end{equation}
Indeed, for $d=p-\unif$, $\lVert d\rVert_2^2\leq\lVert d\rVert_1$
and $\lVert d\rVert_1\leq2\lVert d\rVert_2$, and Jensen's
inequality gives the upper bound after taking expectations.

The proof of the lower bound is the part that permits parameters
arbitrarily close to one. The degree-independent upper bound follows
from a cubic change of variable in the binary moment-saving inequality.
Section~\ref{sec:examples}
explains the weight in \eqref{eq:potential}, treats zero channels, and
shows why local finiteness cannot simply be omitted.
